%%%%%%%%%%%%%%%%%%%%%%%%%%%%%%%%%%%%%%%%%%%%%%%%%%%%%%%%%%%%%%%%%%%%%
%% This is a (brief) model paper using the achemso class
%% The document class accepts keyval options, which should include
%% the target journal and optionally the macuscript tye
%%%%%%%%%%%%%%%%%%%%%%%%%%%%%%%%%%%%%%%%%%%%%%%%%%%%%%%%%%%%%%%%%%%%%
\documentclass[journal=jacsat,manuscript=article]{achemso}

%%%%%%%%%%%%%%%%%%%%%%%%%%%%%%%%%%%%%%%%%%%%%%%%%%%%%%%%%%%%%%%%%%%%%
%% Place any additional packages needed here.  Only include packages
%% which are essential, to avoid problems later.
%%%%%%%%%%%%%%%%%%%%%%%%%%%%%%%%%%%%%%%%%%%%%%%%%%%%%%%%%%%%%%%%%%%%%
\usepackage[version=3]{mhchem} % Formula subscripts using \ce{}
\usepackage[normalem]{ulem}
\usepackage{xcolor}
\usepackage{cleveref}
\usepackage{graphicx}
\usepackage{natbib}
\usepackage{enumitem}
\makeatletter
\setlength\acs@tocentry@width{8.25cm}
\setlength\acs@tocentry@height{4.45cm}
\makeatother

%%%%%%%%%%%%%%%%%%%%%%%%%%%%%%%%%%%%%%%%%%%%%%%%%%%%%%%%%%%%%%%%%%%%%
%% If issues arise when submitting your manuscript, you may want to
%% un-comment the next line.  This provides information on the
%% version of every file you have used.
%%%%%%%%%%%%%%%%%%%%%%%%%%%%%%%%%%%%%%%%%%%%%%%%%%%%%%%%%%%%%%%%%%%%%
%%\listfiles

%%%%%%%%%%%%%%%%%%%%%%%%%%%%%%%%%%%%%%%%%%%%%%%%%%%%%%%%%%%%%%%%%%%%%
%% Place any additional macros here.  Please use \newcommand* where
%% possible, and avoid layout changing macros (which are not used
%% when typesetting).
%%%%%%%%%%%%%%%%%%%%%%%%%%%%%%%%%%%%%%%%%%%%%%%%%%%%%%%%%%%%%%%%%%%%%
%\newcommand*{\mycommand}[1]{\texttt{\emph{#1}}}

%%%%%%%%%%%%%%%%%%%%%%%%%%%%%%%%%%%%%%%%%%%%%%%%%%%%%%%%%%%%%%%%%%%%%
%% Meta-data block
%% ---------------
%% Each author should be given as a separate \author command.
%%
%% Corresponding authors should have an e-mail given after the author
%% name as an \email command.
%%
%% The affiliation of authors is given after the authors; each
%% \affiliation command applies to all preceding authors not already
%% assigned an affiliation.
%%
%% The affiliation takes an option argument for the short name.  This
%% will typically be something like "University of Somewhere".
%%
%% The \altaffiliation macro should be used for new address, etc.
%%%%%%%%%%%%%%%%%%%%%%%%%%%%%%%%%%%%%%%%%%%%%%%%%%%%%%%%%%%%%%%%%%%%%
\author{Ryosuke ISHIZUKA}
\affiliation[ZKA]
{Zkanics F.P.S., SIDE-6 Senriyama-West, Suita, Osaka, 565-0851, Japan}
\email{ryo.ishizuka@gmail.com}

%%%%%%%%%%%%%%%%%%%%%%%%%%%%%%%%%%%%%%%%%%%%%%%%%%%%%%%%%%%%%%%%%%%%%
%% The document title should be given as usual
%% A short title can be given as a *suggestion* for running headers.
%%%%%%%%%%%%%%%%%%%%%%%%%%%%%%%%%%%%%%%%%%%%%%%%%%%%%%%%%%%%%%%%%%%%%
\title[\texttt{achemso} demonstration]
{Li/Na Ions Desolvation in Solvate Ionic Liquids:
Molecular Insights From Non-Polarizable and Machine-Learning Force Fields}

\begin{document}
%%%%%%%%%%%%%%%%%%%%%%%%%%%%%%%%%%%%%%%%%%%%%%%%%%%%%%%%%%%%%%%%%%%%%
%% The manuscript does not need to include \maketitle, which is
%% executed automatically.  The document should begin with an
%% abstract, if appropriate.  If one is given and should not be, the
%% contents will be gobbled.
%%%%%%%%%%%%%%%%%%%%%%%%%%%%%%%%%%%%%%%%%%%%%%%%%%%%%%%%%%%%%%%%%%%%%
\begin{abstract}
%Diffusion of lithium ions in Li-glyme solvate ionic liquids
%and high-concentrate electrolyte were investigated through 
%universal machine learning interatomic potential (MLIP).
%%
%Structural properties.
%%
\end{abstract}

%TOC Graphic.
%\begin{figure}
%\begin{center}
%\includegraphics{TOC}
%\end{center}
%\end{figure}
%Implicit solvation model, Brownian dynamics, Salt effect, XRISM, GB/SA

%%%%%%%%%%%%%%%%%%%%%%%%%%%%%%%%%%%%%%%%%%%%%%%%%%%%%%%%%%%%%%%%%%%%%
%% Start the main part of the manuscript here.
%%%%%%%%%%%%%%%%%%%%%%%%%%%%%%%%%%%%%%%%%%%%%%%%%%%%%%%%%%%%%%%%%%%%%
\clearpage
\section{INTRODUCTION}
Over the past decade, machine-learning force fields (MLFFs) have 
emerged as a transformative approach in computational chemistry.
In these force fields, interatomic interactions are represented by artificial 
neural networks trained on structural data obtained 
from \textit{ab initio} molecular dynamics (AIMD) simulations.
MLFF-based simulations can therefore achieve AIMD-level 
accuracy for many materials.
This success, however, has prompted a reconsideration of 
the validity of the traditionally employed 
non-polarizable force fields.
In particular, the simulation of ionic liquids (ILs) 
is a prototypical case in which non-polarizable 
force fields have been extensively employed.
ILs are molten salts consisting of organic cations 
and organic/inorganic anions.
By varying the combinations of cations and anions,
their electrochemical properties, such as 
ionic conductivity, can be tuned.
ILs have therefore attracted great attention as ``designer solvents''
in electrochemistry and have also been a long-standing subject 
in the community of molecular simulations.
One of the widely used non-polarizable force fields
for ILs is the OPLS-based model developed by Lopes and Padua.
Simulations based on their force fields
can reproduce the fundamental static properties
of ILs with acceptable accuracy.
However, the simulated diffusion of ions 
is much slower than observed experimentally,
which primarily results from the inaccurate modeling 
of electrostatic interactions between ions.
Because the surrounding molecules strongly affect 
the electronic structure of ions, the ionic charge in ILs
(often referred to as effective charge) cannot be 
considered fixed and instead adopts non-integer values.
To account for the fluctuation and transfer of the ionic charge,
considerable effort has been devoted to more realistic modeling 
of electrostatic interactions.
The charge-scaled force field, in which all atomic partial 
charges of ions are uniformly scaled by a non-integer factor, 
has been widely employed as the simplest approach.
This charge scaling effectively weakens the ion-ion interactions,
thereby accelerating ion diffusion.
However, since the scaling factor is empirically determined
to reproduce the experimental properties of ILs,
reliable experimental data are required.
Furthermore, simulations using this force field provide 
limited molecular-level insights due to the inherent lack
of physical justification for the scaling factor.
To overcome this limitation, Ishizuka and Matubayasi 
developed the MD/DFT self-consistent scheme to determine 
the effective charge within the quantum mechanical framework~\cite{Ishizuka2016, Ishizuka2017}.
They successfully determined the effective charge
via the population analysis of the electron density 
of ILs without relying on any experimental data.
In contrast to MLFFs, these methods are based on non-polarizable force fields, 
so the instantaneous forces acting on each atom are not consistent
with the Hellmann-Feynman forces for the corresponding atomic configurations.
Nevertheless, it remains unclear why refinements of
non-polarizable force fields are so successful 
in reproducing both the static and dynamic properties of ILs.
To clarify this, we perform a comparative study of non-polarizable
and machine-learning force fields through simulations of 
${\rm Li}^{+}$ and ${\rm Na}^{+}$ desolvation 
in glyme-based solvate ILs. 

The glyme-based ionic liquids are characterized by their unique 
solvation structures, where a lithium ion is coordinated by a single glyme
molecule to form a stable solvate complex, as represented in Figure~\ref{fig:glyme_ILs}.
%%%%
Recent experimental studies suggest that the several tri-glyme is.
The binding strength between ions and 
%%%%
In the present work, therefore, our interest lies not in the effectiveness 
and performance of this ionic liquid as an electrolyte, 
but in its use as a benchmark for machine-learning and 
non-polarizable force fields.
%%%%
%%%%

\clearpage

\begin{figure}
\begin{center}
\includegraphics{./glyme_ILs}
\caption{ 
Schematic illustrations of (a) a lithium ion solvated by triglyme,
(b) the anions considered in this work, and (c) glyme-based solvate ionic liquids,
exemplified by $[{\rm Li}({\rm G}_{3})][{\rm NO}_{3}]$.
}
\label{fig:glyme_ILs}
\end{center}
\end{figure}

\clearpage

\begin{figure}
\begin{center}
\includegraphics[width=\linewidth]{./fcorr_vdwd3}
\caption{ 
Comparison of atomic forces calculated with PBE and PBE-D3 for 
(a) $[{\rm Li}({\rm G3})][{\rm NO}_{3}]$,
(b) $[{\rm Li}({\rm G3})][{\rm OTf}]$,
(c) $[{\rm Na}({\rm G3})][{\rm NO}_{3}]$,
and (d) $[{\rm Na}({\rm G3})][{\rm OTf}]$.
Forces from PBE are shown on the horizontal axis, 
while those from PBE-D3 are shown on the vertical axis.
}
\label{fig:fcorr_vdwd3}
\end{center}
\end{figure}

%%%%
%Ab initio MD (AIMD) has been performed for tetra glyme ionic liquids~\cite{Sun2018}
%%%%
%Machine-learning interatomic potentials (MLIPs).
%Review of ML \cite{Noe2020, Martin-Barrios2024}.
%Machine-learning force fields (MLFFs).
%CHGnet developed by Zhong et al. \cite{CHGNET}.
%M3Gnet \cite{Chen2022}
%Graph neural network (GNN) based machine-learning force fields.
%Molecular dynamics (MD) simulation has been a powerful tool for the study of
%the structure, dynamics, and function of proteins in atomic detail.
%In general, the result of the MD simulation is quite sensitive to 
%the accuracy of the interatomic interactions.
%The classical force field for ionic liquid 
%%%%%
%%%%%
%$[{\rm Li}({\rm G3})][{\rm NO}_{3}]$ and $[{\rm Li}({\rm G3})][{\rm OTf}]$ 
%are evaluated by experiments \cite{Ueno2012}. 
%%%%%
%

\section{Many-body DPD simulation}

\begin{eqnarray}
{\bf F}_{ij}^{\rm C} &= A_{ij}\omega_{\rm C}(r_{ij}){\hat {\bf r}}_{ij}
+ B_{ij}({\bar\rho}_{i} + {\bar\rho}_{j) \omega_{\rm d}(r_{ij}) 
\label{eq:MDPD_CONS}
\end{eqnarray}




%%%%%%%%%%%%%%%%%%%%%%%%%%%%%%%%%%%%%%%%%%%%%%%%%%%%%%%%%%%%%%%%%%%%%
%% The "Acknowledgement" section can be given in all manuscript
%% classes.  Rather than use \section, an appropriate macro is
%% provided that will always work.
%%%%%%%%%%%%%%%%%%%%%%%%%%%%%%%%%%%%%%%%%%%%%%%%%%%%%%%%%%%%%%%%%%%%%
\acknowledgement
The author gratefully acknowledges the support 
of the Zkanics F.P.S. \\ 
(https://www.linkedin.com/in/zkanics-side6).

%%%%%%%%%%%%%%%%%%%%%%%%%%%%%%%%%%%%%%%%%%%%%%%%%%%%%%%%%%%%%%%%%%%%%
%% The same is true for Supporting Information, which should use the
%% \suppinfo macro.
%%%%%%%%%%%%%%%%%%%%%%%%%%%%%%%%%%%%%%%%%%%%%%%%%%%%%%%%%%%%%%%%%%%%%
%\suppinfo
%
%The entire \textsf{achemso} bundle is generated by running
%\texttt{achemso.dtx} through \TeX. Running \LaTeX\ on the same file
%will generate the general documentation for both the class and
%package files.

%%%%%%%%%%%%%%%%%%%%%%%%%%%%%%%%%%%%%%%%%%%%%%%%%%%%%%%%%%%%%%%%%%%%%
%% The appropriate \bibliography command should be placed here.
%% Notice that the class file automatically sets \bibliographystyle
%% and also names the section correctly.
%%%%%%%%%%%%%%%%%%%%%%%%%%%%%%%%%%%%%%%%%%%%%%%%%%%%%%%%%%%%%%%%%%%%%
\bibliography{reference}

\end{document}
